% Abstract
\begin{abstract}
Current AI alignment methods optimize models to match human preferences through Reinforcement Learning from Human Feedback (RLHF), treating alignment as unidirectional---adjusting AI policy until humans are satisfied. This framing ignores that conversations are bidirectional: users also learn to communicate effectively with AI over time. We propose behavioral coupling as a metric for bidirectional alignment, quantified through the correlation between two observable signals: user learning rate (reformulation slope) and AI responsiveness (query-response diversity correlation). Analyzing 169,352 multi-turn conversations from the HH-RLHF dataset, we find that users reformulate queries less frequently over turns (mean slope = $-0.021$, $p=0.012$, Cohen's $d=-0.246$), indicating learning of AI capabilities, and AI response diversity correlates with query diversity ($r=0.396$, $p<0.001$), indicating behavioral responsiveness. Correlation strengthens with conversation length ($r=0.04$ for 2--3 turns $\rightarrow$ $r=0.19$ for 6+ turns), consistent with alignment emerging through extended interaction rather than being a static policy property. While effect sizes are smaller than initially hypothesized, these findings validate the theoretical foundation of bidirectional alignment and demonstrate that co-adaptation signals can be extracted from standard conversation logs without additional human evaluation. This work shifts alignment measurement from static policy quality to conversation dynamics, enabling new training objectives that optimize for co-adaptation rather than unidirectional preference matching.
\end{abstract}
